\documentclass[singlespacing,draft,titlepage=firstiscover,a4paper,11pt,twoside,numbers=noenddot]{scrbook}
\usepackage[onehalfspacing]{setspace}
\AfterTOCHead{\singlespacing}
\KOMAoptions{DIV=last}
\interfootnotelinepenalty=10000
\usepackage{amsmath}
\usepackage{amssymb}
\usepackage{amsthm}
\usepackage{tensor}
\usepackage[utf8]{inputenc}
\usepackage[T1]{fontenc}
\usepackage{lmodern}
\usepackage[ngerman,english]{babel}
\usepackage{csquotes}
\usepackage[backend=biber,style=numeric]{biblatex}
\addbibresource{thesis.bib}
\usepackage[final,babel=true]{microtype}
\usepackage{multicol}
\usepackage{multirow}
\usepackage{bigdelim}
\usepackage[final]{graphicx}
\usepackage[shortlabels]{enumitem}
\usepackage{makecell}
\usepackage{booktabs}
\usepackage{longtable}
\usepackage{siunitx}
\usepackage{nomencl}
\makenomenclature
\usepackage[svgnames]{xcolor}
\usepackage{pgf}
\graphicspath{{fig/}}
\renewcommand{\pgfimage}[2][]{\includegraphics[#1]{#2}}
\usepackage{subcaption}
\captionsetup[subfigure]{justification=centering}
\usepackage{gnuplot-lua-tikz}
\usepackage{ifdraft}
\usetikzlibrary{external}
\usepackage{algorithmic}
\usepackage[final]{listings}
\usepackage{eqparbox}
\usepackage[final]{hyperref}
\usepackage{orcidlink}
\usepackage[obeyFinal]{todonotes}
\usepackage{s4probe}

%% unit definitions
% -----------------------------------------
\sisetup{per-mode=symbol, range-phrase=--, range-units=single, tight-spacing=true, text-series-to-math=true, propagate-math-font=true}
\DeclareSIUnit\Byte{\text{Byte}} % use full name 'Byte' instead of 'B'
\DeclareSIUnit\flop{\text{Flop}}
\DeclareSIUnit\cycle{\text{cycle}}
\DeclareSIUnit\instruction{\text{instruction}}
% helpful abbreviations
\DeclareSIUnit\kB{\kilo\byte}
\DeclareSIUnit\MB{\mega\byte}
\DeclareSIUnit\GB{\giga\byte}
\DeclareSIUnit\TB{\tera\byte}
\DeclareSIUnit\PB{\peta\byte}
\DeclareSIUnit\GBs{\giga\byte\per\second}
\DeclareSIUnit\GFs{\giga\flop\per\second}
\DeclareSIUnit\TBs{\tera\byte\per\second}
\DeclareSIUnit\TFs{\tera\flop\per\second}
% For the ERA5 data
\DeclareSIUnit\hPa{\hecto\pascal}
% -----------------------------------------

%% This code creates groups for the nomenclature
% -----------------------------------------
\usepackage{etoolbox}
\renewcommand\nomgroup[1]{%
\item[\bfseries%
    \ifstrequal{#1}{A}{Number sets}{}%
    \ifstrequal{#1}{B}{Scalars}{}%
    \ifstrequal{#1}{C}{Vectors}{}%
    \ifstrequal{#1}{D}{Matrices}{}%
    \ifstrequal{#1}{E}{Tensors}{}%
    \ifstrequal{#1}{F}{Tensor operators}{}%
    \ifstrequal{#1}{P}{Physical quantities}{}%
    ]%
}
% -----------------------------------------

%% two-column nomenclature
% -----------------------------------------
\renewcommand{\nompreamble}{\begin{multicols}{2}}
\renewcommand{\nompostamble}{\end{multicols}}
% -----------------------------------------

%% addition theorem environments
% -----------------------------------------
\newtheorem{remark}{Remark}[chapter]
\newtheorem{theorem}{Theorem}[chapter]
\newtheorem{proposition}{Proposition}[chapter]
\newtheorem{disclaimer}{Disclaimer}[chapter]
\newtheorem{supposition}{Supposition}[chapter]
\newtheorem{corollary}{Corollary}[chapter]
\newtheorem{lemma}{Lemma}[chapter]
% -----------------------------------------

%% autoref names (use non-capital first letter)
% -----------------------------------------
\addto\extrasenglish{\def\chapterautorefname{Chapter}}
\addto\extrasenglish{\def\sectionautorefname{Section}}
\addto\extrasenglish{\def\subsectionautorefname{Section}}
\addto\extrasenglish{\def\figureautorefname{Fig.}}
\addto\extrasenglish{\def\tableautorefname{Table}}
\addto\extrasenglish{\def\algorithmautorefname{Alg.}}
\addto\extrasenglish{\def\remarkautorefname{Remark}}
\addto\extrasenglish{\def\corollaryautorefname{Corollary}}
\addto\extrasenglish{\def\lemmaautorefname{Lemma}}
\addto\extrasenglish{\def\propositionautorefname{Proposition}}
\newcommand*{\subsecref}[1]{\hyperref[{#1}]{\textbf{\nameref*{#1}} on page~\pageref*{#1}}}
% -----------------------------------------

%% Macros to abbreviate common math symbols and operators
% -----------------------------------------
\newcommand{\R}{\mathbb{R}}
\renewcommand{\Rn}{\mathbb{R}^n}
\newcommand{\Rm}{\mathbb{R}^m}
\newcommand{\Rnm}{\mathbb{R}^{n\times m}}
\newcommand{\Rnn}{\mathbb{R}^{n\times n}}
\newcommand{\C}{\mathbb{C}}
\newcommand{\Cn}{\mathbb{C}^n}
\newcommand{\Cnm}{\mathbb{C}^{n\times m}}
\newcommand{\N}{\mathbb{N}}
\newcommand{\TT}{\text{TT}}
\DeclareMathOperator*{\argmin}{argmin}
\DeclareMathOperator*{\argmax}{argmax}
\DeclareMathOperator{\rank}{rank}
\DeclareMathOperator{\ranks}{ranks}
\DeclareMathOperator{\diag}{diag}
\DeclareMathOperator{\reshape}{reshape}
\DeclareMathOperator{\unfold}{unfold}
\DeclareMathOperator{\trunc}{trunc}
\newcommand{\contract}{\;\Join\;}
\newcommand{\contractij}[2]{\;\tensor[_{#1}]{\Join}{_{#2}}\;}
% -----------------------------------------

%% Colors (picked from modern gnuplot)
% -----------------------------------------
\definecolor{gnuplot1}{RGB}{148,0,211}
\definecolor{gnuplot2}{RGB}{0,158,115}
\definecolor{gnuplot3}{RGB}{86,180,233}
\definecolor{gnuplot4}{RGB}{230,159,0}
% -----------------------------------------

%% Setup for algorithms
% -----------------------------------------
\renewcommand{\algorithmicrequire}{\textbf{Input:}}
\renewcommand{\algorithmicensure}{\textbf{Output:}}
\renewcommand{\algorithmiccomment}[1]{\hfill| \eqparbox{COMMENT\thealgorithm}{\textsl{\color{darkgray}#1}}\ignorespaces}
%% instead of \usepackage{algorithm}
%\usepackage{algorithm}
\DeclareNewTOC[
    type=algorithm,
    float,floattype=4,floatpos=h,counterwithin=chapter,
    atbegin={%
        {\noindent\rule[\ht\strutbox]{\linewidth}{\heavyrulewidth}\vspace{-1.3\baselineskip}}%
        %\KOMAoptions{captions=heading}%
        \xpretocmd\algorithmic
            {\noindent\rule[\ht\strutbox]{\linewidth}{\lightrulewidth}\vspace{-2\baselineskip}}
            {}{\PatchFailed}%
        \xapptocmd\endalgorithmic
            {\noindent\rule[\ht\strutbox]{\linewidth}{\heavyrulewidth}\vspace{-\baselineskip}}
            {}{\PatchFailed}%
    }
]{loalg}
% -----------------------------------------

%% Setup for source listings
% -----------------------------------------
% when several languages are needed, this only works by defining own styles!
\lstset{
    language=C++,
    frame=tb,
    xleftmargin=.1\textwidth, xrightmargin=.1\textwidth,
    numbers=left,
    basicstyle=\linespread{1.0}\scriptsize\ttfamily,
    keywordstyle=\color{DarkBlue}\ttfamily,
    stringstyle=\color{DarkRed}\ttfamily,
    commentstyle=\color{DarkGreen}\ttfamily,
    morekeywords={__shared__},
    morecomment=[l][\color{DarkMagenta}]{\#}
}
%\lstset{language=C++}
%\lstset{deletedirectives={\#pragma},
%}
%    morekeywords={__shared__},
% -----------------------------------------


%% Faster compiling with pre-generated tikz pictures
% -----------------------------------------
\ifdraft{%
\tikzexternalize[shell escape=-enable-write18,prefix=fig_tmp/]
% except todo notes - they don't work with tikzexternaliz
\NewCommandCopy\origtodo\todo
\RenewDocumentCommand\todo{O{}m}{%
  \tikzexternaldisable\origtodo[#1]{#2}\tikzexternalenable}
}{%
% Nothing to do
}
% -----------------------------------------

\begin{document}
\frontmatter
\tableofcontents
\listoffigures
\listoftables
\listoftodos
\printnomenclature
\mainmatter
\chapter{Equations and floats}\label{ch:one}
\section{Theory}\label{sec:theory}
Text with a citation \cite{knuth84}\footnote{A footnote.} and \SI{3}{\GBs}.
\nomenclature[A]{$\R$}{real numbers}\todo{check this}
\begin{equation}\label{eq:one} x \contract y \in \Rn \end{equation}
\begin{align} a &= b \\ c &= d \end{align}
\begin{theorem}\label{thm:one} A theorem. \end{theorem}
\begin{figure}[htb]\centering\includegraphics[width=3cm]{a}\caption{A figure}\label{fig:one}\end{figure}
\begin{table}[htb]\centering\caption{A table}\label{tab:one}
\begin{tabular}{lll}\toprule \multirow{2}{*}{A} & \makecell{B\\C} & D \\ & E & F \\\bottomrule\end{tabular}\end{table}
\section*{Unnumbered}
More text.

\chapter{Tables and algorithms}\label{ch:two}
See \autoref{eq:one} and \nameref{sec:theory}.\nomenclature[B]{$x$}{a scalar}
\begin{longtable}{ll}\caption{Long table}\label{tab:long}\\\toprule A & B\\\midrule\endhead
row 1 & value 1\\
row 2 & value 2\\
row 3 & value 3\\
row 4 & value 4\\
row 5 & value 5\\
row 6 & value 6\\
row 7 & value 7\\
row 8 & value 8\\
row 9 & value 9\\
row 10 & value 10\\
row 11 & value 11\\
row 12 & value 12\\
row 13 & value 13\\
row 14 & value 14\\
row 15 & value 15\\
row 16 & value 16\\
row 17 & value 17\\
row 18 & value 18\\
row 19 & value 19\\
row 20 & value 20\\
row 21 & value 21\\
row 22 & value 22\\
row 23 & value 23\\
row 24 & value 24\\
row 25 & value 25\\
row 26 & value 26\\
row 27 & value 27\\
row 28 & value 28\\
row 29 & value 29\\
row 30 & value 30\\
row 31 & value 31\\
row 32 & value 32\\
row 33 & value 33\\
row 34 & value 34\\
row 35 & value 35\\
row 36 & value 36\\
row 37 & value 37\\
row 38 & value 38\\
row 39 & value 39\\
row 40 & value 40\\
row 41 & value 41\\
row 42 & value 42\\
row 43 & value 43\\
row 44 & value 44\\
row 45 & value 45\\
row 46 & value 46\\
row 47 & value 47\\
row 48 & value 48\\
row 49 & value 49\\
row 50 & value 50\\
row 51 & value 51\\
row 52 & value 52\\
row 53 & value 53\\
row 54 & value 54\\
row 55 & value 55\\
row 56 & value 56\\
row 57 & value 57\\
row 58 & value 58\\
row 59 & value 59\\
row 60 & value 60\\
row 61 & value 61\\
row 62 & value 62\\
row 63 & value 63\\
row 64 & value 64\\
row 65 & value 65\\
row 66 & value 66\\
row 67 & value 67\\
row 68 & value 68\\
row 69 & value 69\\
row 70 & value 70\\
\bottomrule\end{longtable}
\begin{algorithm}\caption{An algorithm}\label{alg:one}\begin{algorithmic}\REQUIRE $x$ \COMMENT{input}\STATE $y \gets x$ \COMMENT{assign a longer comment}\ENSURE $y$\end{algorithmic}\end{algorithm}
\begin{lstlisting}
int main() { return 0; }
\end{lstlisting}
\begin{enumerate}\item one\item two\end{enumerate}
\begin{figure}[htb]\centering\begin{subfigure}{.4\textwidth}\centering\includegraphics[width=2cm]{a}\caption{Left}\end{subfigure}\begin{subfigure}{.4\textwidth}\centering\includegraphics[width=2cm]{a}\caption{Right}\end{subfigure}\caption{Subfigures}\end{figure}

\chapter{Pictures}\label{ch:three}
\begin{tikzpicture}\draw (0,0) -- (2,1) node[right]{tikz};\end{tikzpicture}
\section*{Starred}
\subsecref{sec:theory}, \autoref{tab:long}, \autoref{alg:one}, \cite{lamport94}.
\foreignlanguage{ngerman}{Ein deutscher Satz.}\todo{another note}

\appendix
\chapter{Appendix}\label{ch:app}
Appendix text, \autoref{ch:one}, page~\pageref{fig:one}.

\printbibliography
\end{document}
